\documentclass[12pt,a4paper]{article}
\usepackage[utf8]{vietnam}
\usepackage[margin=2.5cm]{geometry}
\usepackage{tcolorbox}
\usepackage{fancyhdr}

\pagestyle{empty}

\begin{document}

\begin{center}
    \textbf{\Large CỘNG HÒA XÃ HỘI CHỦ NGHĨA VIỆT NAM}\\
    \textbf{\large Độc lập - Tự do - Hạnh phúc}\\
    \rule{8cm}{0.5pt}\\[1cm]
    
    \textbf{\Huge DI CHÚC}\\
    \textit{(Văn bản viết tay gốc thu giữ trong hộp sắt)}
\end{center}

\vspace{1cm}

\begin{large}
Tôi tên là: \textbf{Nguyễn Văn Thành}, sinh năm 1940.\\
Chủ sở hữu hợp pháp căn nhà số 14, Đường Bờ Sông.

\vspace{0.5cm}

Nay tôi còn hoàn toàn minh mẫn lập di chúc này để lại toàn bộ quyền sử dụng đất và tài sản gắn liền với đất tại căn nhà số 14 Đường Bờ Sông (thuộc \textbf{Mã thửa địa chính khóa 2021-BS14}) cho cháu nội tôi là \textbf{Nguyễn Văn Khang} toàn quyền sở hữu sau khi tôi qua đời.

Cháu Trần Ngọc Mai không được chia phần nhà này.

\vspace{1.5cm}

\begin{flushright}
    \textit{Bờ Sông, ngày 15 tháng 04 năm 2018}\\
    \textbf{Người lập di chúc}\\
    \vspace{1.5cm}
    \textbf{Nguyễn Văn Thành} (Đã ký)
\end{flushright}
\end{large}

\vspace{2cm}

\begin{tcolorbox}[colback=red!5!white,colframe=red!75!black,title=\textbf{MANH MỐI BẰNG CHỨNG GIẢ MẠO (EUREKA CLUE)}]
\begin{itemize}
    \item Ngày ký di chúc ghi: \textbf{15/04/2018}.
    \item Thuật ngữ ghi trong văn bản: \textbf{"Mã thửa địa chính khóa 2021-BS14"} — Hệ thống mã này chỉ được Bộ Tài nguyên áp dụng từ năm 2021!
    \item \textbf{Kết luận:} Tờ di chúc bị Khang viết đè chèn nội dung giả mạo sau năm 2021.
\end{itemize}
\end{tcolorbox}

\end{document}
